\chapter{Discussion and Limitations}
\label{ch:discussion}

\section{Interpretation of the Main Result}

The headline AUROC of 0.94 must be read together with the composition of the cohort. Participants with MCI converted 25 times as often as cognitively normal participants, and the diagnosis alone reached an AUROC of 0.84. Much of the overall discrimination therefore comes from separating these groups, which a clinician already does. The more informative figures are the within-group results: about 0.84 among participants with MCI, and about 0.88 among cognitively normal participants, the latter based on few events. Removing all information about current cognition lowered the AUROC from 0.94 to 0.75. The model is therefore best described as estimating the short-term risk of progression given the current clinical stage, not as detecting future dementia in people without any signs of impairment.

Including CDR-SB and diagnosis as inputs is not leakage in the technical sense: both are measured at the index visit, and the outcome is a diagnosis made at a later visit. It does mean that the clinical label used as outcome and some of the inputs come from the same assessment process. Results without these inputs are reported for this reason.

\section{The Cardiovascular Question}

The central hypothesis of the original project was that cardiovascular information would improve prediction. In this cohort it added about 0.001 AUROC once current cognition was known, the survival analysis gave the same result, and the hazard ratios for hypertension, diabetes and stroke were close to one after adjustment. This is consistent with, rather than contrary to, the epidemiological literature. Most of the evidence concerns midlife exposure or exposure accumulated over decades \cite{sierra2020,wang2026}, whereas NACC participants enter at a median age of 71 and their history is recorded largely as yes/no answers with single blood-pressure readings. Damage that vascular disease has already caused is likely to be reflected in the cognitive scores, which the model sees directly. The finding should therefore be stated as limited incremental predictive value in this cohort and setting, not as evidence that vascular health is unimportant for the brain.

\section{Comparison with Related Studies}

The finding that cardiovascular information adds little once other information is known agrees with Aschwanden et al., who found cardiovascular factors less important than expected when evaluated jointly with other predictors \cite{aschwanden2020}. It does not contradict the association studies \cite{jia2023,song2020}, which ask whether cardiovascular risk is related to later cognition rather than whether it improves a model that already measures current cognition. Our AUROC is higher than that of risk models without cognitive inputs, such as the UK Biobank model of Stern and Linial (0.773) \cite{stern2026}, because our model includes current cognition and clinical stage; without them our model scored 0.75, in the same range. Like Yi et al.\ \cite{yi2023}, we found CDR-SB to be the most important feature, although their task was diagnosis rather than prediction of future conversion.

\section{Model Complexity}

Logistic regression, random forest, XGBoost, a multilayer perceptron and an FT-Transformer performed within 0.007 AUROC of one another, with overlapping intervals. This agrees with comparisons on other tabular data \cite{gorishniy2021}. With a few dozen informative features, most of them strongly related to the outcome, there is little structure left for a more flexible model to exploit. The modality-encoder fusion network proposed at the start of the project gave the same result, and its attention weights were nearly uniform, so they could not serve as explanations. The deep-learning idea that did prove useful was modality dropout, which addresses a practical problem (missing assessments) rather than raising discrimination.

\section{PET}

Amyloid status was strongly associated with progression in participants with MCI, in line with the biological definition of the disease \cite{jack2024}, but no clear gain was shown when PET was added to the clinical model. The likely reason is limited follow-up: most scans are recent, and only a few dozen scanned participants with a known three-year outcome converted. The PET image experiment showed that the raw images can be processed into measurements that agree closely with the PET core's, but with five converters among the participants with images, image-based prediction of dementia was not attempted.

\section{Limitations}

\begin{enumerate}
\item \textbf{Cohort.} NACC is a research cohort recruited through memory centres and volunteers; participants are mostly White and highly educated. Absolute risks and performance may differ in other populations.
\item \textbf{Validation.} The held-out centres belong to the same study, use the same forms and were collected in the same period. This is site-level validation, not external validation on an independent cohort, and no temporal validation was performed.
\item \textbf{Censoring.} About half of the eligible participants were excluded because their three-year outcome was unknown. Participants with more severe impairment return less often, so the analysed cohort is somewhat healthier than all enrolled participants. The survival analysis that includes censored participants reached the same conclusions.
\item \textbf{Outcome.} The outcome is a clinical diagnosis, which can be revised; 7.7\% of converters were later recorded as not demented. Excluding them did not change the results.
\item \textbf{Cognitively normal participants.} Only 24 test and 37 held-out converters were cognitively normal at the index visit, so performance in this group is uncertain.
\item \textbf{Calibration.} Risk was slightly over-predicted at the held-out centres and by about eight percentage points for participants with MCI younger than 65.
\item \textbf{Demographic inputs.} Race, ethnicity and living situation improved calibration across racial groups, most likely because of differences in recruitment, but they change individual estimates noticeably. Whether such inputs belong in a clinical tool is a question of policy as well as accuracy.
\item \textbf{Test reuse.} The same test set was used for several experiments. The held-out centres provide a second evaluation set, and the results agree.
\item \textbf{PET images.} Registration was performed without MRI, tau images were not quantified, and Grad-CAM maps were not informative.
\item \textbf{Evidence assistant.} The library contains eleven papers and the test questions were written by the team.
\end{enumerate}

\section{Ethical Considerations}

The data were used under the NACC data use agreement. No participant-level information leaves the analysis machine, and every shared output is checked automatically for identifiers. The prototype is labelled as a research tool, does not diagnose, and does not recommend treatment: the data record what happened to participants and cannot show what would have happened under a different treatment. The evidence assistant answers only from published sources and declines to give advice about an individual.
